\documentclass[10pt,a4paper]{amsart}
\usepackage[margin=19mm]{geometry}
\usepackage{amsmath,amssymb,mathtools}
\usepackage[scr=rsfs,cal=euler]{mathalfa}
\usepackage{xfrac}
\usepackage[unicode,hidelinks,bookmarks=false]{hyperref}
\newcommand{\ie}{i.e.\ }
\newcommand{\cf}{cf.\ }
\newcommand{\resp}{respectively\ }
\newcommand{\Subsection}{\subsection}
\providecommand{\nobreakdash}{\nobreak\hbox{-}}
\setlength{\emergencystretch}{2em}
\multlinegap0pt
\usepackage{mathtools}
\usepackage{amssymb}
\usepackage[shortlabels]{enumitem}
\usepackage{float}
\newtheorem{lemma}{Lemma}
\DeclareMathOperator{\Tr}{Tr}
\newcommand{\dt}{\ {\rm d} t }
\newcommand{\dx}{\ {\rm d} {x}}
\title{A mathematical study of an elastic-viscous-plastic sea-ice model with the Kelvin-Voigt rheology}
\author{Daniel W. Boutros, Xin Liu, Marita Thomas, Edriss S. Titi}
\date{Source: arXiv:2604.26295v1 (CC BY 4.0)}
\begin{document}
\maketitle

\noindent\emph{Source and licence.} A mathematical study of an elastic-viscous-plastic sea-ice model with the Kelvin-Voigt rheology. Version arXiv:2604.26295v1 (CC BY 4.0), distributed by arXiv under the Creative Commons Attribution 4.0 International licence, http://creativecommons.org/licenses/by/4.0/, as recorded in the arXiv licence metadata for this exact version and shown on the abstract page https://arxiv.org/abs/2604.26295v1. Full licence notice: \texttt{licenses/arxiv-2604.26295-CC-BY-4.0.txt}.\par\smallskip

\noindent\textit{Editorial scope.} Counted unit: the article's lemma Lpestimatelemma (source0.tex lines 280--289). The article's complete original proof of it is delivered with it (source0.tex lines 290--330, one \texttt{proof} environment of 41 lines). No mathematics is written, repaired or re-derived: the statement and the proof are the article's own lines, in the article's own notation, and no retained line is substituted or re-flowed.  Retained uncounted as the source-local context the proof needs: lines 202--207; lines 211--216.  Nothing was omitted from any retained span, so the package carries no omission record.  The retained lines cite 1 works, whose entries are transcribed verbatim from the article's own bibliography source in the e-print and delivered with short editorial labels recorded in the item's reference mapping.  No retained line contains a figure environment, an image-inclusion command or drawing code, so no image is delivered and the package declares no asset dependency.  Editorial formatting only, in addition: an AMS \texttt{amsart} preamble that restates the article's own declarations and macros used by the retained lines, namely the environments \texttt{lemma} on separate counters, together with the standard packages those lines need.  The retained environments print in this document as Lemma 1 in order of appearance, whereas the article prints the same items with the numbering of its own full document.  The complete native source of the article as distributed in the arXiv e-print for this exact version is bundled unchanged as \texttt{sources/199381/source0.tex}.
\begin{align}
&\partial_t (u - \alpha^2 \Delta u) + u \cdot \nabla u  = \nabla \cdot \sigma + \mathcal{T}_a + \mathcal{T}_w + \Omega u^\perp - g \nabla H_0, \label{advection1epsilon} \\
&\dfrac{1}{\mathcal E} \partial_t \sigma + \dfrac{4 \mathcal{D}_\epsilon}{P} (\sigma - \frac{1}{2} \Tr \sigma \mathbb I_2) + \dfrac{\mathcal{D}_\epsilon}{2P} \Tr \sigma \mathbb I_2 + \dfrac{\mathcal{D}_\epsilon}{2} \mathbb I_2  = D(u). \label{advection2epsilon}
\\
&u \lvert_{t = 0} = u_0, \quad \sigma \lvert_{t = 0} = \sigma_0. \label{advection3epsilon}
\end{align}

\begin{align}
&\partial_t (u - \alpha^2 \Delta u) = \nabla \cdot \sigma + \mathcal{T}_a + \mathcal{T}_w + \Omega u^\perp - g \nabla H_0, \label{kelvinvoigt1epsilon} \\
&\dfrac{1}{\mathcal E} \partial_t \sigma + \dfrac{4 \mathcal{D}_\epsilon}{P} (\sigma - \frac{1}{2} \Tr \sigma \mathbb I_2) + \dfrac{\mathcal{D}_\epsilon}{2P} \Tr \sigma \mathbb I_2 + \dfrac{\mathcal{D}_\epsilon}{2} \mathbb I_2  = D(u), \label{kelvinvoigt2epsilon}
	\\
&u \lvert_{t = 0} = u_0, \quad \sigma \lvert_{t = 0} = \sigma_0, \label{kelvinvoigt3epsilon}
\end{align}

\begin{lemma} \label{Lpestimatelemma}
Let $u \in C([0,T];H^1 (\mathbb{T}^2))$ and $\sigma \in C([0,T];H^1_w (\mathbb{T}^2)) \cap C([0,T]; L^2 (\mathbb{T}^2))$, with $\sigma_0 \in L^\infty (\mathbb{T}^2) \cap H^1 (\mathbb{T}^2)$, be two functions which satisfy the following equation (where $P > 0$ is a given constant)
\begin{equation}
\dfrac{1}{\mathcal E} \partial_t \sigma + \dfrac{4 \mathcal{D}_\epsilon}{P} (\sigma - \frac{1}{2} \Tr \sigma \mathbb I_2) + \dfrac{\mathcal{D}_\epsilon}{2P} \Tr \sigma \mathbb I_2 + \dfrac{\mathcal{D}_\epsilon}{2} \mathbb I_2  = D(u),
\end{equation}
which is the same as \eqref{advection2epsilon} and \eqref{kelvinvoigt2epsilon}. Then it holds that $\sigma \in L^\infty ((0,T); L^\infty (\mathbb{T}^2))$ and we have
\begin{equation}
\lVert \sigma (\cdot, t) \rVert_{L^\infty} \leq \lVert \sigma_0 \rVert_{L^\infty} + 2 P.
\end{equation}
\end{lemma}

\begin{proof}
We first introduce the new unknown
\begin{equation} \label{newunkown}
\tau \coloneqq \sigma + \frac{P}{2} \mathbb I_2,
\end{equation}
which satisfies the following equation
\begin{equation} \label{modifiedstressequation}
\dfrac{1}{\mathcal E} \partial_t \tau + \dfrac{4 \mathcal{D}_\epsilon}{P} (\tau - \frac{1}{2} \Tr \tau \mathbb I_2) + \dfrac{\mathcal{D}_\epsilon}{2P} \Tr \tau \mathbb I_2 = D(u).
\end{equation}
From the regularity of $u$ and $\tau$ it follows that $\partial_t \tau \in C([0,T]; L^{2 - \gamma} (\mathbb{T}^2))$ for any $\gamma \in (0,1)$. Therefore, for every $p > 3$, we can take the duality pairing of equation \eqref{modifiedstressequation} with $P^{-p+1} \tau \lvert \tau \rvert^{p-2}$, which gives
\begin{align*}
&\frac{P}{p \mathcal{E}} \frac{\textrm{d}}{\dt} \lVert \tau / P \rVert_{L^p}^p + \int_{\mathbb{T}^2} \bigg[ 4 \mathcal{D}_\epsilon \left( \frac{1}{P} \right)^p \bigg\lvert \tau - \frac{1}{2} \Tr \tau \mathbb I_2 \bigg\rvert^2 \lvert \tau \rvert^{p-2} + \dfrac{\mathcal{D}_\epsilon}{2} \left( \frac{1}{P} \right)^p \lvert \Tr \tau \rvert^2 \lvert \tau \rvert^{p-2} \bigg] \dx \\
&\leq \int_{\mathbb{T}^2} P^{-p+1} \mathcal{D}_\epsilon \lvert \tau \rvert^{p-1} \dx = \int_{\mathbb{T}^2} \mathcal{D}_\epsilon^{1/p} \cdot P^{-(p-1)} \mathcal{D}_\epsilon^{(p-1)/p} \lvert \tau \rvert^{p-1} \dx \\
&\leq \int_{\mathbb{T}^2} \bigg[\frac{1}{p} \mathcal{D}_\epsilon + \frac{p-1}{p} \mathcal{D}_\epsilon \left( \frac{1}{P} \right)^p \lvert \tau \rvert^p \bigg] \dx.
\end{align*}
Now, we recall the property
\begin{equation*}
\lvert \tau \rvert^2 = \bigg\lvert \tau - \frac{1}{2} \Tr \tau \mathbb I_2 \bigg\rvert^2 + \frac{1}{2} \lvert \Tr \tau \rvert^2.
\end{equation*}
Therefore, we can deduce the following estimate
\begin{align*}
\frac{P}{p \mathcal{E}} \frac{\textrm{d}}{\dt} \lVert \tau / P \rVert_{L^p}^p &+ \int_{\mathbb{T}^2} \bigg[ 3 \mathcal{D}_\epsilon \left( \frac{1}{P} \right)^p \bigg\lvert \tau - \frac{1}{2} \Tr \tau \mathbb I_2 \bigg\rvert^2 \lvert \tau \rvert^{p-2} + \dfrac{\mathcal{D}_\epsilon}{p} \left( \frac{1}{P} \right)^p \lvert \tau \rvert^p \bigg] \dx \leq \frac{1}{p} \int_{\mathbb{T}^2} \mathcal{D}_\epsilon \dx.
\end{align*}
Dropping the coercive terms (the positive terms on the left-hand side) and integrating in time leads to
\begin{align}
\dfrac{1}{\mathcal E} \lVert (\tau / P) (\cdot, t) \rVert_{L^p}^p \leq \dfrac{1}{\mathcal E} \lVert \tau_0 / P \rVert_{L^p}^p + \frac{1}{P} \int_0^t \int_{\mathbb{T}^2} \mathcal{D}_\epsilon \dx \dt' \leq \dfrac{1}{\mathcal E} \lVert \tau_0 / P \rVert_{L^p}^p + \frac{1}{P} \int_0^t (\lVert \nabla u \rVert_{L^2} + \epsilon) \dt' \label{integratedLpestimate}.
\end{align}
Now, taking the $p$-th root of equation \eqref{integratedLpestimate} and using the subadditivity property of the $p$-th root (i.e. $(x + y)^{1/p} \leq x^{1/p} + y^{1/p}$ for $x,y \geq 0$, $ p \geq 1 $) we have
\begin{equation} \label{integratedLpestimate2}
\dfrac{1}{\mathcal{E}^{1/p}} \lVert (\tau / P) (\cdot, t) \rVert_{L^p} \leq \dfrac{1}{\mathcal{E}^{1/p}} \lVert \tau_0 / P \rVert_{L^p} + \left( \frac{1}{P} \int_0^t (\lVert \nabla u \rVert_{L^2} + \epsilon) \dt' \right)^{1/p},
\end{equation}
which holds for every $p > 3$. Next, we recall the following result for finite-measure spaces: If $f \in L^p (\mathbb{T}^2)$ for any $p \in [p_0,\infty)$, for some $p_0 \geq 1$, and $\lVert f \rVert_{L^p} \leq K$ (for a constant $K$ independent of $p$), then $f \in L^\infty (\mathbb{T}^2)$ and $\lVert f \rVert_{L^\infty} \leq K$. Moreover, we have
\begin{equation} \label{Linftylimit}
\lVert f \rVert_{L^\infty} = \lim_{p \rightarrow \infty} \lVert f \rVert_{L^p},
\end{equation}
the proof of which can be found for example in \cite[Theorem 2.8]{adams}. 
Now sending $p \rightarrow \infty$ in equation \eqref{integratedLpestimate2} and using property \eqref{Linftylimit} we find
\begin{equation}
\lVert \tau (\cdot, t) \rVert_{L^\infty} \leq \lVert \tau_0 \rVert_{L^\infty} + P.
\end{equation}
\end{proof}

\begin{thebibliography}{99}
\bibitem[adams]{adams} {\sc Adams, R.~A.} \newblock {\em Sobolev Spaces}. \newblock Academic Press, 1975.
\end{thebibliography}
\end{document}
